% FRAME Paper 1 — ICLR draft skeleton (v0.1, pre-Stage-2)
% Compile: requires iclr2026_conference.sty + .bst (or swap the two style
% lines below for \documentclass{article} + natbib to compile standalone).
% Placeholders for Stage-2 results are marked \pending{...} in orange.
\documentclass{article}
\usepackage{iclr2026_conference,times}   % ICLR style; comment out if absent
%\usepackage{times}                      % standalone fallback
\usepackage{natbib}
\usepackage{amsmath,amssymb}
\usepackage{booktabs}
\usepackage{xcolor}
\usepackage{hyperref}
\usepackage{microtype}

\newcommand{\pending}[1]{\textcolor{orange!80!black}{[#1]}}
\newcommand{\kap}{\kappa}
\newcommand{\framework}{\textsc{Frame}}

\title{The Directive Reaches the Rationale, Not the Decision:\\
Reliability-Gated, Pre-Registered Measurement of\\
Political-Framing Behavior in Frontier Language Models}

\author{Anonymous Authors\\
\texttt{anonymous@example.com}}

\begin{document}
\maketitle

\begin{abstract}
Neutrality directives---system-prompt instructions such as \emph{``do not
editorialize; represent perspectives proportionally''}---are the de facto
industry remedy for political bias in deployed language models. We test
whether they do what they appear to do. In a pre-registered, two-stage study
of two frontier models (Claude Sonnet 4.5, GPT-4.1) operating on a balanced
corpus of 200 real news articles stratified by expert-rated political lean,
we separate what a directive changes about a model's \emph{decisions} (which
bias it detects; how it classifies an article's lean) from what it changes
about the model's \emph{stated rationale}. Our design is built around a
methodological commitment we argue the LLM-judge literature needs:
\textbf{measurement precedes findings}. Every judge instrument had to pass a
cross-family inter-rater reliability gate ($\kap \ge 0.40$, threshold and
fallback committed in writing before validation data existed, one attempt,
no re-rolling). Two of our three instruments \emph{failed}---one judge
scored 89\% of items at the scale ceiling; another never used one of two
labels across 60 items---and the pilot's most striking apparent finding
(models re-use loaded source framing in 55--64\% of explanations) proved to
be largely an artifact of judge error, collapsing to $\sim$10\% under a
corrected rubric. The surviving instrument ($\kap = 0.606$) measures the
political direction of framing substitution. With it, we test the
pre-registered prediction that directives leave decisions unchanged
(equivalence tests on detection counts and lean labels) while stated
rationales drift directionally---with right-leaning sources sanitized toward
the center more than left-leaning sources. Stage-2 results:
\pending{H26/H28/H29 outcomes}. Beyond the findings, which are snapshots of
two models at one moment, we release the durable artifact: a deterministic,
refreshable corpus recipe, gated instruments, and a pre-registration
protocol that lets any future model generation be audited identically.
\end{abstract}

\section{Introduction}

When a language model is asked to analyze a politically charged news
article, it produces two things: a \emph{decision} (an answer---which spans
are biased, what the article's lean is) and a \emph{rationale} (a stated
explanation of that answer). Almost all governance of political behavior in
deployed models operates on the assumption that these move together: add a
neutrality directive to the system prompt, observe better-sounding
explanations, and conclude the model has become more even-handed. This paper
tests that assumption directly, and pre-registers the prediction that it is
false in a specific, directional way.

Our central design principle is a dissociation test. If a directive
genuinely alters processing, its effects should appear in the decisions
themselves. If instead the directive is absorbed as a \emph{presentation
instruction}, decisions will be unchanged while rationales are rewritten to
\emph{perform} compliance. The distinction matters for AI safety practice
well beyond the political domain: it determines whether behavioral audits
and self-report audits of the same system can be expected to agree. A model
whose explanations respond to instructions that its decisions ignore will
pass every audit that reads its explanations and fail every audit that
measures its behavior---or, more dangerously, the reverse.

Testing this required confronting a second problem, which became a
co-equal contribution of the paper. At the scale of thousands of judgments,
the only affordable instruments are LLM judges, and the field's default
practice is to write a plausible rubric, spot-check a few outputs, and
trust the resulting numbers. We instead adopted the standard that
psychometrics applies to any new instrument: \textbf{demonstrate
reliability before collecting confirmatory data, under acceptance criteria
committed in advance}. Concretely, every instrument was scored
independently by two frontier judges from different model families
(cross-family pairing, to break same-family circularity), on held-out pilot
data excluded from all confirmatory analysis, against a pre-committed gate
(Cohen's $\kap \ge 0.40$ plus non-degeneracy conditions), with a
pre-registered fallback and a one-shot rule forbidding iterative
re-anchoring until the gate passes.

The gates did real work. Of three instruments that all looked reasonable on
paper---and whose pilot outputs all \emph{looked} sensible to
inspection---two failed calibration in ways that would have silently
corrupted the study (\S\ref{sec:gates}). Most consequentially, the pilot's
most publishable number, an apparent 55--64\% rate at which models absorb
loaded source framing into their own explanatory voice, was traced by
forensic audit to judge error: under a corrected, evidence-anchored rubric
the rate collapsed to roughly 10\%. We report this trajectory in full not
as a confession but as a finding: plausible-looking LLM-judge instruments
fail cross-family calibration at high rates, the failures are systematic
rather than noisy, and unn gated instruments convert judge artifacts into
publishable effects.

\textbf{Contributions.}
\begin{enumerate}
  \item \textbf{A dissociation result on instruction-faithfulness}
  (pre-registered; Stage-2 outcomes \pending{confirmed/disconfirmed}):
  neutrality directives leave frontier models' decisions statistically
  unchanged (equivalence on detection counts, $|\Delta| < 2.0$ per article;
  label stability $\kap \ge 0.85$) while their stated rationales shift, with
  a directional signature---right-leaning sources' framing is stripped or
  re-coded toward the center \pending{X}$\times$ more often than
  left-leaning sources' (pilot prior: $2$--$3\times$, both judge families
  agreeing on sign independently).
  \item \textbf{A measurement-integrity protocol with empirical teeth}:
  pre-committed reliability gates for LLM-judge instruments, under which 2
  of 3 plausible instruments failed; a demonstration that the gates caught
  a large artifact before it became a headline finding; and a
  characterization of the failure modes (anchor divergence, base-rate
  degeneracy, off-definition abstention) with the rubric-design features
  that distinguished the surviving instrument (categorical labels,
  quoted-evidence requirements, explicit abstention semantics).
  \item \textbf{A durable, refreshable audit capability}: a deterministic
  corpus-generation recipe over expert-rated news (contamination-resistant
  by regeneration), gated instruments, analysis code, and a pre-registration
  with externally timestamped amendments---packaged so that the study can be
  re-run on any future model generation for $\sim$\$150.
\end{enumerate}

We emphasize the asymmetry between our findings and our instrument. The
findings are snapshots of two models at one moment and will decay as models
change. The question the instrument asks---\emph{do instructions reach the
decision, or only the story about the decision?}---must be re-answered by
every model generation. The durable contribution is the validated capacity
to ask it.

\section{Motivation and Related Work}
\label{sec:motivation}

\textbf{From opinion surveys to behavior.} A substantial literature
documents left-of-center tendencies in LLM outputs, largely by
administering political-orientation questionnaires or eliciting stances on
contested issues \citep[e.g.,][]{rozado2024politics}. Recent critiques
observe that such measurements cannot distinguish a model's
\emph{viewpoint preferences} from failures of \emph{epistemic integrity},
and that directional asymmetry is not per se evidence of bias when inputs
themselves are asymmetric \citep{rozado2026vpei}. We adopt this distinction
rather than dispute it, and operationalize what it leaves unmeasured: not
what models \emph{say they believe}, but what they \emph{do} to political
framing when performing realistic editorial tasks---detection,
summarization, classification---on real news, with the article's outlet and
headline withheld to prevent source-prior shortcuts. Our pre-registered
justified-asymmetry companion analysis (\S\ref{sec:setup}) directly
addresses the strongest objection this literature would raise against our
headline test: per-stratum intensity balance is measured (objective loaded-lexicon
density; expert rating magnitudes; the models' own pre-directive detection
counts), with the interpretation rule fixed before data collection.

\textbf{Faithfulness of stated reasoning.} \citet{turpin2023language}
showed that chain-of-thought explanations can omit the features that
actually drive predictions. Our design asks the complementary question: not
whether explanations \emph{omit} true causes, but whether an explicit
instruction changes the explanation \emph{without} changing the decision it
purports to explain---and whether that rewriting has a political direction.
The generation-order arms (rationale emitted before vs.\ after the label)
probe the same post-hoc hypothesis from a second angle.

\textbf{LLM judges and their validation.} LLM-as-judge evaluation is known
to exhibit position, verbosity, and self-preference biases
\citep{zheng2023judging}, and agreement with human raters varies widely
across tasks. What the literature largely lacks is \emph{consequence}:
reliability is measured, lamented, and the judges are used anyway. Our
protocol imports the missing step from psychometrics and clinical-trial
practice \citep{landis1977measurement}: instruments earn confirmatory
status by passing acceptance criteria fixed in advance, or they are
demoted---a rule we enforced against our own study twice
(\S\ref{sec:gates}). We are not aware of prior work in which pre-committed
reliability gates with registered fallbacks were applied to LLM-judge
instruments before confirmatory data collection; we document what such
gates catch.

\textbf{Construct validity in NLP evaluation.} Our approach follows calls
to treat evaluation instruments as measurement devices requiring validity
arguments rather than as ground truth \citep{jacobs2021measurement}. The
gate-failure post-mortems (\S\ref{sec:gates}) supply concrete instances:
identical rubrics resolved in opposite, internally consistent ways by two
frontier judges; Cohen's $\kap$ collapsing under extreme base rates while
raw agreement stands at 87\% \citep{gwet2008computing}.

\section{Experimental Setup}
\label{sec:setup}

\textbf{Corpus.} $N{=}200$ contemporary U.S.\ news articles on a balanced
$5\times8$ grid: five AllSides-anchored lean classes (Left, Lean~Left,
Center, Lean~Right, Right; 40 each) crossed with eight topical themes
(25 each), 400--1500 words, per-outlet caps within lean, deterministic
selection from a 10{,}000-article enriched pool by a released curation
script. Articles were boundary- and interior-cleaned; an identity audit
(id$\leftrightarrow$url$\leftrightarrow$text) was performed after six
mismatched rows were discovered in the pilot pool (all excluded from
analysis; none present in the final corpus). Models receive
\emph{article text only}---no headline, outlet, URL, or date---so that
classification cannot reduce to outlet priors.

\textbf{Targets and tasks.} Two frontier models (Claude Sonnet 4.5;
GPT-4.1) perform three tasks: \textbf{Eval~A}, span-level bias detection
with per-detection explanations; \textbf{Eval~B}, summarization with key
facts and represented perspectives; \textbf{Eval~C}, five-class lean
classification with a mandatory free-text rationale. Structured JSON
outputs; one rollout per cell.

\textbf{Conditions.} Each task runs under prompt arms that differ by
exactly one block (verified by programmatic diff): \emph{baseline} (task
only), \emph{ablation} (adds bias-type schema and vocabulary; no
directive), \emph{reframing} (ablation $+$ the neutrality directive),
\emph{full} (ablation $+$ an attribution rule), and, for Evals A and C, a
\emph{reframing-CoT} arm that reorders generation so the rationale precedes
the decision. The registered directive contrast is
reframing~vs.~ablation.

\textbf{Instruments and gates.}
\label{sec:instruments}
Judge instruments are scored independently by a cross-family pair (Claude
Sonnet 4.6; GPT-5) on identical inputs. Three were built:
\textbf{VAR} (voice adoption: does an explanation \emph{describe} loaded
source framing or \emph{inherit} it?), \textbf{FDC} (1--7 frame-distance
coding on attribution and schema axes), and \textbf{Directional RD}
(categorical: does an output's framing substitute \emph{leftward},
\emph{rightward}, neither, or is there no signal---plus a
strip/amplify mechanism, with quoted-evidence requirements). Each faced a
pre-committed acceptance gate on held-out pilot data (excluded from all
confirmatory analysis): $\kap \ge 0.40$ with non-degeneracy conditions, one
attempt, registered fallback (demotion to exploratory, per-judge
reporting). Human calibration on a stratified subsample
(\pending{$n$, coders, agreement}) arbitrates all demoted constructs.

\textbf{Confirmatory hypotheses} (Benjamini--Hochberg FDR, $q{=}.05$,
family of three after gate outcomes; full decision rules in the
pre-registration):
\textbf{H26}: on Eval-C rationales, $P(\text{left-substitution} \mid
\text{Right-lean source}) > P(\text{right-substitution} \mid
\text{Left-lean source})$; primary test judge-pooled with a per-judge
sign-consistency requirement (a significant pooled result with discordant
per-judge signs counts as non-confirmation).
\textbf{H28}: equivalence (TOST, $|\Delta| < 2.0$ detections/article) of
Eval-A detection counts, reframing vs.\ ablation.
\textbf{H29}: stability of Eval-C lean labels across the directive
contrast (bootstrap $\kap \ge 0.85$).
H28 and H29 require no judge; H26 uses the sole gate-passing instrument.

\textbf{Pre-registration and integrity.} All hypotheses, thresholds,
fallbacks, exclusions, and amendments are timestamped in a public git
history; every instrument revision was committed \emph{before} the data
that evaluated it existed. The two demotions below were executed exactly
as registered.

\section{Findings}

\subsection{Finding 1: Most plausible judge instruments fail calibration---and gates catch what inspection does not}
\label{sec:gates}

This finding is complete (it precedes Stage-2 data by design).
Table~\ref{tab:gates} summarizes.

\begin{table}[h]
\centering
\small
\begin{tabular}{lllll}
\toprule
Instrument & Scale type & Gate statistic & Outcome & Failure mode \\
\midrule
FDC & 1--7 Likert, 2 axes & qw-$\kap = 0.109$ & \textbf{fail} &
anchor divergence: one judge 89\% at ceiling \\
& & & & ($\sigma{=}0.47$); 2.5-pt systematic offset after repair \\
VAR & categorical (3) & $\kap = 0.000$ & \textbf{fail} &
zero-variance rater (0/60 one label); raw agreement \\
& & & & 86.7\%, AC$_1{=}0.848$: base-rate degeneracy \\
Directional RD & categorical (4) $+$ mech. & $\kap = 0.606$ &
\textbf{pass} & --- (from $\kap{=}0.233$ pre-anchoring) \\
\bottomrule
\end{tabular}
\caption{Pre-committed reliability gates ($\kap \ge 0.40$, one shot,
registered fallbacks). Cross-family judge pairs; held-out pilot items.}
\label{tab:gates}
\end{table}

Three observations. \emph{First}, both failures were systematic, not
noisy: each judge applied an articulable, internally consistent rule that
the rubric failed to adjudicate (e.g., whether generic source-mention
constitutes attribution). \emph{Second}, the pilot's most striking
apparent result---55--64\% of explanations inheriting loaded source
framing---was substantially a judge artifact: forensic audit found 47\% of
``inheriting'' verdicts cited phrases absent from the source; under the
corrected rubric the rate fell to $\sim$10\%. Ungated, this artifact was
publishable. \emph{Third}, the surviving instrument differed by design
features, not luck: categorical labels rather than Likert anchors,
quoted-evidence requirements for non-default verdicts, and explicit
abstention semantics. Reliability rose from $\kap = 0.233$ to $0.606$ after
these were added, while the same repair strategy could not save the Likert
instrument.

\subsection{Finding 2: Decisions are unchanged by the directive (H28, H29)}
\pending{Stage-2: TOST outcome, equivalence bounds, per-target detail;
label-stability $\kap$ with bootstrap CI. Pilot priors indicate detection
counts and lean labels move little across arms.}

\subsection{Finding 3: Rationales drift, directionally (H26)}
\pending{Stage-2: pooled estimate with cluster-robust CI; per-judge sign
consistency; no-signal coverage by arm and stratum. Pilot prior (excluded
corpus, both judges independently): right-lean sources' framing
re-coded toward center 2--3$\times$ more often than left-lean sources'.}

\subsection{Finding 4: The asymmetry is not explained by input intensity (D-JA companion)}
Corpus-side indicators, computed pre-collection: objective loaded-lexicon
density is balanced across Left and Right strata (1.73 vs.\ 1.93 per 1k
words; $d{=}0.10$, $p{=}.55$); expert rating magnitudes are imbalanced
\emph{against} the predicted direction (Left articles rated more intense),
making proportionate-response accounts of the predicted asymmetry
conservative. \pending{Stage-2: models' own baseline detection counts per
stratum; intensity-normalized asymmetry if any indicator is imbalanced.}

\section{Discussion and Limitations}

\textbf{What the dissociation means if confirmed.} A model whose rationale
responds to a neutrality instruction that its decisions ignore is not
``more neutral''; it is more fluent at narrating neutrality. For
practitioners this cuts in a specific direction: system-prompt directives
are presentation-layer controls, and any governance regime that audits
models by reading their explanations will systematically mis-estimate
behavioral change. Behavioral audits---counting decisions under controlled
contrasts---are not optional rigor; they measure a different quantity.

\textbf{What the gate failures mean regardless of Stage-2 outcomes.} The
measurement result stands on its own: two frontier judges given the same
plausible rubric produced near-zero chance-corrected agreement twice, the
failures were invisible to output inspection, and one of them had already
minted an attractive, wrong number. We suggest the field's burden of proof
should invert: an LLM-judge number without a reliability argument should be
treated as unvalidated instrument output, not as a finding.

\textbf{Limitations.} (1) \emph{Scope}: two closed-weight models, U.S.\
politics, English news; the pattern's generality is untested (the pipeline
re-runs on additional targets for marginal cost). (2) \emph{No mechanism}:
we establish behavioral dissociation, not its cause; generation-order arms
probe but cannot settle it. (3) \emph{Judge mediation}: H26 rides on one
validated instrument; cross-family sign agreement and human calibration
(\pending{results}) bound, but do not eliminate, shared-judge-bias risk.
(4) \emph{Gate statistics}: $\kap$ is base-rate sensitive; our own VAR
post-mortem illustrates the $\kap$ paradox (high raw agreement, degenerate
marginals), which is why gates included non-degeneracy conditions and why
demotions route to human arbitration rather than deletion. (5)
\emph{Construct floor}: under the corrected rubric, explanation-level
framing inheritance appears rare; whether that reflects model behavior or
rubric strictness awaits human coding. (6) \emph{One rollout per cell};
sampling variance is handled across articles, not within.

\section{Conclusion}

We set out to measure whether neutrality directives change what frontier
models do with politically framed news, and found we first had to build
instruments that could be trusted to measure anything at all. The
resulting study contributes a pre-registered dissociation test between
decisions and rationales (\pending{outcome}), a directional
characterization of framing substitution on real news under a
justified-asymmetry discipline (\pending{outcome}), and---independent of
either---a demonstrated protocol under which evaluation instruments earn
trust before they are believed. The findings are dated the moment new
models ship. The capability to re-derive them, on a regenerated corpus
with re-gated instruments for roughly the cost of a conference
registration, is the contribution we expect to last.

\subsubsection*{Ethics Statement}
This study evaluates model behavior on politically contested material; we
have tried to hold it to the epistemic standard we advocate. The research
question is directional symmetry, not the vindication of any political
position; the corpus is balanced by construction against expert ratings;
the loaded-language lexicon pairs left- and right-coded terms; judge
prompts were audited for direction-priming and example asymmetries, with
two such asymmetries found and corrected pre-data. We adopt the
distinction between viewpoint preference and epistemic failure and
explicitly do not equate measured asymmetry with ``bias'': our
pre-registered companion analysis tests whether asymmetric treatment
tracks asymmetric inputs, and our claims are conditioned on its outcome.
Misuse risks run in two directions---overreading a directional default as
partisan intent, and weaponizing ``neutrality'' audits to punish accurate
asymmetry \citep{rozado2026vpei}; our reporting rules (per-judge
estimates, intensity normalization, equivalence tests reported alongside
directional ones) are designed against both. News text is used under
research norms with per-outlet caps; article text is redistributed only as
required for reproduction, with source URLs and the deterministic curation
script provided in lieu of bulk redistribution where licensing requires.
Planned human annotation involves paid coders rating published news text
(no personal data, no vulnerable populations); compensation and consent
procedures follow institutional guidelines. All API costs
($\sim$\$200) and the full audit trail of instrument failures are
disclosed to make the study's error-correction process, not only its
conclusions, inspectable.

\subsubsection*{Reproducibility Statement}
The pre-registration (hypotheses, thresholds, gates, fallbacks, and every
amendment) is timestamped in a public git history, with instrument
revisions committed before the data that evaluated them existed. We
release: the deterministic corpus curation script and corpus manifest, all
judge prompts (versioned, with the full revision history), the runner and
analysis code, gate validation data, batch manifests, and per-call cost
accounting. A regeneration of the full study on new articles or new
target models requires one configuration change and $\sim$\$150 in batched
API calls. \pending{Repository and OSF links after de-anonymization.}

\bibliography{references}
\bibliographystyle{iclr2026_conference}

% --- references.bib (starter; verify entries before submission) ---------
% @article{rozado2024politics,
%   title={The political preferences of LLMs},
%   author={Rozado, David}, journal={PLOS ONE}, year={2024}}
% @misc{rozado2026vpei,
%   title={Have We Been Measuring AI Political Bias Wrong?},
%   author={Rozado, David and others}, year={2026},
%   howpublished={Substack essay, VPEI framework}}
% @inproceedings{turpin2023language,
%   title={Language Models Don't Always Say What They Think: Unfaithful
%          Explanations in Chain-of-Thought Prompting},
%   author={Turpin, Miles and Michael, Julian and Perez, Ethan and
%           Bowman, Samuel R.}, booktitle={NeurIPS}, year={2023}}
% @inproceedings{zheng2023judging,
%   title={Judging LLM-as-a-Judge with MT-Bench and Chatbot Arena},
%   author={Zheng, Lianmin and others}, booktitle={NeurIPS Datasets and
%   Benchmarks}, year={2023}}
% @inproceedings{jacobs2021measurement,
%   title={Measurement and Fairness},
%   author={Jacobs, Abigail Z. and Wallach, Hanna},
%   booktitle={FAccT}, year={2021}}
% @article{landis1977measurement,
%   title={The measurement of observer agreement for categorical data},
%   author={Landis, J. Richard and Koch, Gary G.},
%   journal={Biometrics}, year={1977}}
% @article{gwet2008computing,
%   title={Computing inter-rater reliability and its variance in the
%          presence of high agreement},
%   author={Gwet, Kilem L.}, journal={British Journal of Mathematical and
%   Statistical Psychology}, year={2008}}

\end{document}
